\paragraph{Two-regime identities.}
Let $E$ be a two-state continuous-time Markov chain and let the conditional
genealogy generator in regime $e$ be $G_n(p_e)^T$. The joint generator is
\[
 \widetilde A_n=R\otimes I_{d_n}+\operatorname{blockdiag}
 (G_n(p_L)^T,G_n(p_H)^T).
\]
For four lineages, $d_4=19$ and therefore
$\dim(\widetilde A_4)=38$.

\begin{proposition}[Equal regime frequencies]
If $p_L=p_H=p$, then the marginal quartet law of the two-regime model equals
the fixed-frequency law at $p$, for every conservative regime generator $R$,
initial regime distribution, and set of branch durations, provided all other
genealogical parameters are equal between regimes.
\end{proposition}
\begin{proof}
The two diagonal genealogy blocks are identical, so the joint semigroup is the
Kronecker sum of $R$ and the common genealogy generator. The regime factor
preserves total mass and disappears upon marginalization.
\end{proof}

\begin{proposition}[Taxon-exchange symmetry]
Suppose a permutation of taxa exchanges $Q_2$ and $Q_3$ and leaves the complete
parameterized species-tree, terminal sampling, branch parameters, and
inheritance kernels invariant. Then the unconditional law satisfies
$q_2=q_3$.
\end{proposition}
\begin{proof}
Relabeling gives a measure-preserving bijection of regime and genealogy paths
and exchanges the two absorbing topology labels. The two probabilities are
therefore equal.
\end{proof}

The numerical search accompanying this proposition gives a deterministic
witness with $q_3>q_1$ and $q_2\ne q_3$. This is an existence witness for the
implemented model, not a universal analytic anomaly criterion. Full-simplex
coverage is not established.
